\begin{center}
\begin{minipage}{\linewidth}
\centering
\resizebox{0.96\linewidth}{!}{%
\begin{tikzpicture}[
  >=Latex,
  every node/.style={font=\small},
  ws/.style={very thick,draw=Azul},
  wu/.style={very thick,draw=Rojo},
  flow/.style={-{Latex[length=2.3mm]},very thick}
]
  \draw[rounded corners=6pt,thick,GrisOscuro] (0.10,0.10) rectangle (6.35,5.75);
  \node[font=\bfseries] at (3.22,5.42) {Retorno homoclínico en \(\Sigma\)};

  \coordinate (p) at (2.35,2.10);
  \draw[ws]
    (2.18,0.38) .. controls (1.85,1.35) and (2.55,3.25) .. (2.25,5.10);
  \draw[wu]
    (0.45,2.36) .. controls (1.18,2.05) and (1.70,2.02) .. (p)
    .. controls (4.15,1.95) and (5.55,2.55) .. (5.20,3.72)
    .. controls (4.90,4.78) and (3.35,4.65) .. (2.30,4.25)
    .. controls (1.30,3.88) and (0.80,3.32) .. (0.55,2.80);
  \fill[GrisOscuro] (p) circle (2.7pt);
  \node[below right] at (p) {$p$};
  \coordinate (q) at (2.30,4.25);
  \fill[GrisOscuro] (q) circle (2.7pt);
  \node[left=5pt] at (q) {$q$};
  \node[Azul,fill=white,inner sep=2pt] at (1.50,4.72) {$W^s(p)$};
  \node[Rojo,fill=white,inner sep=2pt] at (4.62,4.45) {$W^u(p)$};
  \draw[flow,Rojo] (3.10,2.08) -- (3.82,2.18);
  \draw[flow,Rojo] (4.80,3.28) -- (4.62,3.90);
  \draw[flow,Azul] (2.22,3.25) -- (2.27,2.62);
  \draw[flow,Azul] (2.25,4.82) -- (2.28,4.43);
  \node[align=center,text=GrisOscuro,font=\scriptsize] at (3.25,0.72)
    {$q\in W^s(p)\cap W^u(p)$\\
      $T_qW^s+T_qW^u=T_q\Sigma$};

  \draw[-{Latex[length=2.7mm]},very thick,GrisOscuro]
    (6.70,2.90) -- node[above,align=center] {Smale--\\Birkhoff} (8.05,2.90);

  \draw[rounded corners=6pt,thick,GrisOscuro] (8.35,0.10) rectangle (13.95,5.75);
  \node[font=\bfseries] at (11.15,5.42) {Estirar, contraer y plegar};
  \draw[line width=16pt,Azul!60,rounded corners=12pt]
    (9.25,0.42) -- (9.25,5.00) -- (12.45,5.00) -- (12.45,0.42);
  \draw[line width=7pt,white,rounded corners=9pt]
    (9.25,0.42) -- (9.25,5.00) -- (12.45,5.00) -- (12.45,0.42);
  \draw[very thick,GrisOscuro] (8.78,0.85) rectangle (12.92,4.38);
  \node[anchor=north west] at (8.88,4.28) {$R$};
  \node[font=\bfseries] at (9.25,2.35) {$R_0$};
  \node[font=\bfseries] at (12.45,2.35) {$R_1$};
  \draw[-{Latex[length=2.3mm]},thick,GrisOscuro]
    (10.05,1.35) -- (10.05,0.72);
  \draw[-{Latex[length=2.3mm]},thick,GrisOscuro]
    (11.65,1.35) -- (11.65,0.72);
  \node[draw=GrisOscuro,rounded corners,fill=white,inner sep=3pt]
    at (11.15,0.42)
    {$\cdots010011\cdots\in\{0,1\}^{\mathbb Z}$};
\end{tikzpicture}%
}

\smallskip
\begin{minipage}{0.92\linewidth}
\small\textbf{Lectura pedagógica.}
Para un mapa de retorno, una intersección homoclínica transversal de las
variedades de un punto hiperbólico permite construir tiras que se codifican por
símbolos. La flecha representa una implicación teórica bajo las hipótesis de
Smale--Birkhoff; no afirma que la intersección ya haya sido certificada en Chua
o MAVPD.
\end{minipage}
\end{minipage}
\end{center}
